% Target: rmp-iii-b-self-braiding
% Formula: ring_j wound twice = e^{i phi_j} ring_j — the anyon-j MPO
% Ink: Canonical public kernel: one declared self-crossing string carries eight
%   tensors before resolving to the four-tensor ring.
% Boundary: Both MPO cycles close; one marked j-string tail remains open in each
%   panel.
% Source: paper TN-Review-main.tex line 1662; author source ImagesReview Section
%   III/ImagesReview Section III.tex lines 704-738
% Capabilities: kernel, self-braiding, crossing-order
\tenkzkernel
\tndeclare{species}{operator}{hue=source:red}
\[
  \begin{tenkz}[rows={wire,wire,wire,wire,wire}, cols=5, bonds=none]
    % One closed string follows the outer turn, then the inner turn.  The
    % declared self-crossing is the flattened braid at the source's
    % upper-right seam; the crossing police rejects the panel if that
    % topology drifts.  The closed via ring is twelve invisible junction
    % atoms jd1..jd12 chained by twelve named wires d1..d12, d_k running
    % P_k -> P_{k+1}.  Three joints (P3, P5, P7) continue in a straight
    % line and need no declared crossing; the other nine bend and each
    % carries a nominal cross={over at crossing of ...} declaration at the
    % shared junction point (no visible gap: the "crossing" is the
    % segments' own shared endpoint, not an interior overlap).  The one
    % interior self-crossing survives as a real declared order between d9
    % (the inner top edge) and d12 (the closing diagonal), where the
    % source's flattened braid seam actually falls.
    \tn[at=(1,3), skin=none, name=jd1]{} \tn[at=(1,5), skin=none, name=jd2]{}
    \tn[at=(3,5), skin=none, name=jd3]{} \tn[at=(5,5), skin=none, name=jd4]{}
    \tn[at=(5,3), skin=none, name=jd5]{} \tn[at=(5,1), skin=none, name=jd6]{}
    \tn[at=(3,1), skin=none, name=jd7]{} \tn[at=(1,1), skin=none, name=jd8]{}
    \tn[at=(2,2), skin=none, name=jd9]{} \tn[at=(2,4), skin=none, name=jd10]{}
    \tn[at=(4,4), skin=none, name=jd11]{} \tn[at=(4,2), skin=none, name=jd12]{}
    \tnwire[kind=string, species=operator, name=d1,
      cross={over at crossing of d1 and d2}]{jd1}{jd2}
    \tnwire[kind=string, species=operator, name=d2]{jd2}{jd3}
    \tnwire[kind=string, species=operator, name=d3,
      cross={over at crossing of d3 and d4}]{jd3}{jd4}
    \tnwire[kind=string, species=operator, name=d4]{jd4}{jd5}
    \tnwire[kind=string, species=operator, name=d5,
      cross={over at crossing of d5 and d6}]{jd5}{jd6}
    \tnwire[kind=string, species=operator, name=d6]{jd6}{jd7}
    \tnwire[kind=string, species=operator, name=d7,
      cross={over at crossing of d7 and d8}]{jd7}{jd8}
    \tnwire[kind=string, species=operator, name=d8,
      cross={over at crossing of d8 and d9}]{jd8}{jd9}
    \tnwire[kind=string, species=operator, name=d9,
      cross={over at crossing of d9 and d10,
             under at crossing of d9 and d12}]{jd9}{jd10}
    \tnwire[kind=string, species=operator, name=d10,
      cross={over at crossing of d10 and d11}]{jd10}{jd11}
    \tnwire[kind=string, species=operator, name=d11,
      cross={over at crossing of d11 and d12}]{jd11}{jd12}
    \tnwire[kind=string, species=operator, name=d12,
      cross={over at crossing of d12 and d1}]{jd12}{jd1}
    \tn[skin=box, at=on d1 0.4915]{} \tn[skin=box, at=on d2 0.9661]{}
    \tn[skin=box, at=on d4 0.5636]{} \tn[skin=box, at=on d6 0.0382]{}
    \tn[skin=box, at=on d7 0.6357]{}
    \tn[skin=box, role=marked, name=mkD, label pos=ne, at=on d9 0.4031]{j}
    \tn[skin=box, at=on d11 0.0006]{} \tn[skin=box, at=on d12 0.3005]{}
    \tnwire[kind=string, species=operator, name=twD,
      cross={over at crossing of twD and d12}]{mkD.n}{n outside}
  \end{tenkz}
  \;=\; e^{i\phi_j}\;
  \begin{tenkz}[rows={wire,wire,wire,wire,wire}, cols=5, bonds=none]
    % The closed via ring is six invisible junction atoms js1..js6 chained
    % by six named wires; every joint bends, so every adjacent pair carries
    % a nominal cross={over at crossing of ...} declaration (no visible
    % gap, since the "crossing" is the segments' own shared endpoint).
    \tn[at=(1,3), skin=none, name=js1]{} \tn[at=(2,5), skin=none, name=js2]{}
    \tn[at=(4,5), skin=none, name=js3]{} \tn[at=(5,3), skin=none, name=js4]{}
    \tn[at=(4,1), skin=none, name=js5]{} \tn[at=(2,1), skin=none, name=js6]{}
    \tnwire[kind=string, species=operator, name=s1,
      cross={over at crossing of s1 and s2}]{js1}{js2}
    \tnwire[kind=string, species=operator, name=s2,
      cross={over at crossing of s2 and s3}]{js2}{js3}
    \tnwire[kind=string, species=operator, name=s3,
      cross={over at crossing of s3 and s4}]{js3}{js4}
    \tnwire[kind=string, species=operator, name=s4,
      cross={over at crossing of s4 and s5}]{js4}{js5}
    \tnwire[kind=string, species=operator, name=s5,
      cross={over at crossing of s5 and s6}]{js5}{js6}
    \tnwire[kind=string, species=operator, name=s6,
      cross={over at crossing of s6 and s1}]{js6}{js1}
    \tn[skin=box, at=on s1 0.4631]{}
    \tn[skin=box, role=marked, name=mkS, label pos=ne, at=on s3 0.0159]{j}
    \tn[skin=box, at=on s4 0.4631]{} \tn[skin=box, at=on s6 0.0159]{}
    \tnwire[kind=string, species=operator, name=twS]{mkS.e}{e outside}
  \end{tenkz}
\]
